\subsection{A finite symplectic counterexample in the source chain}

\begin{result}{D.26}{Refuted}{The integral conjugacy assertion in Boggi's Lemma 5.12}
Lemma~5.12 on page~26 of the retrieved version~5 of
\emph{On the procongruence completion of the Teichm\"uller modular group}
asserts that two nonconjugate finite cyclic subgroups of
$\operatorname{Sp}_{2g}(\mathbb Z)$ remain nonconjugate modulo
every integer $m\geq3$.  This assertion already fails for
two order-two subgroups of $\operatorname{Sp}_4(\mathbb Z)$
at $m=3$.  The exact counterexample below therefore also
prevents using the printed proof of Lemma~5.10, and its
subsequent use in Section~5, without an additional repair.
This statement does not assert a counterexample to the
separate profinite twist-normalizer conclusions.
\end{result}
\begin{proof}
Put
\[
 J=\begin{pmatrix}0&1\\-1&0\end{pmatrix},\qquad
 \Omega=\operatorname{diag}(J,J),\qquad
 A=\begin{pmatrix}I_2&0\\0&-I_2\end{pmatrix},\qquad
 B=\begin{pmatrix}0&I_2\\I_2&0\end{pmatrix}.
\]
Both $A$ and $B$ have order two and preserve the integral
symplectic form $\Omega$.  For an integral involution $T$
on $L=\mathbb Z^4$, the index
\[
       [L:L^+(T)+L^-(T)],\qquad
       L^\pm(T)=\ker(T\mp I:L\to L),
\]
is invariant under integral conjugacy: a conjugating
integral automorphism carries the two displayed kernels
onto the corresponding kernels for its conjugate.
For $A$ the two coordinate planes give index $1$.
For $B$ the two eigenlattices are
\[
 L^+(B)=\{(v,v):v\in\mathbb Z^2\},\qquad
 L^-(B)=\{(v,-v):v\in\mathbb Z^2\}.
\]
Their sum consists of the pairs $(x,y)$ with
$x\equiv y\pmod2$, and hence has index $4$.
Thus $A$ and $B$ are not conjugate even in
$\operatorname{GL}_4(\mathbb Z)$.  Their order-two cyclic
subgroups are not conjugate either, since their unique
nonidentity elements would have to be conjugate.

Now take
\[
 Q=\begin{pmatrix}
 1&0&1&0\\
 0&2&0&2\\
 1&0&-1&0\\
 0&2&0&-2
 \end{pmatrix}.
\]
Direct multiplication gives
\[
       BQ=QA,\qquad \det Q=16,\qquad
       Q^{\mathsf t}\Omega Q=4\Omega.
\]
Consequently the reduction $\overline Q$ is symplectic
and invertible modulo $3$, and
$\overline Q\,\overline A\,\overline Q^{-1}=\overline B$.
This is the required conjugacy in
$\operatorname{Sp}_4(\mathbb F_3)$.
All matrix products, determinants and eigenlattice
indices were independently executed and saved in
\path{code/d_symplectic_involutions.py}, with its JSON
certificate and text output.

The exact printed lemma and the paragraph immediately
preceding it were opened and read in the source
\cite[p.~26]{ET_Boggi2014}.  That paragraph explicitly
invokes this lemma to finish Lemma~5.10.  The fixed-locus
argument printed for Lemma~5.12 therefore cannot establish
its claimed implication; the two matrices above give
its precise failed conclusion.  The cutting theorem
used in D.19 and Part~E is Theorem~4.5, whose proof
precedes this argument and was checked separately.
\end{proof}
